% =====================================================================
\section*{How this work was done}
\addcontentsline{toc}{section}{How this work was done}
% =====================================================================

The authors chose the question, directed the work and judged every result. AI models from two developers explored the problem, derived and drafted the arguments, and challenged each other's results: OpenAI's \OPENAIMODEL{} and Anthropic's Claude Fable~5.2 and Claude Opus~5.5. The forgetting theorem and the construction of Theorem~C were each derived by one model and then re-derived refutation-first, step by step, by a model of the other developer, which did not write them. The memory theorem rests on the contour lemma of the forgetting proof, which was cross-checked in this way; its remaining steps were written and checked by models of one developer, the check in a separate run that did not write them, and were checked again by a model of the other developer in the review of version~1.3.0 described below. These reviews found one gap in the parser of \cref{sec:scaffold} (a missing transport of the phase from the anchor to a queried block in the point certificate) and one constant in the memory bound that rested on an unproved property; both were repaired and re-checked. The three-token theorem of \cref{sec:three} was derived by \OPENAIMODEL{} and re-derived refutation-first by Claude Opus~5.5. That review found that the source of the argument used the halting time of $M$ where the halting time of the simulating machine $M^\circ$ is needed; \cref{sec:three} uses the latter. Its one computed lemma was checked by two independent programs with positive controls, and the reviewer re-ran both and wrote a third. The perturbed-kernel corollary of \cref{pf:sec-robust}, added in version~1.2.0, was derived by \OPENAIMODEL{} and re-derived refutation-first, in two rounds, by Claude Opus~5.5. The reconstruction of the manuscript from the proofs was itself read refutation-first by a reviewer of the other developer; for \cref{sec:three}, added in version~1.1.0, and for \cref{pf:sec-robust}, added in version~1.2.0, it was read refutation-first by a model of one developer, in a separate run that did not write it. For version~1.3.0, a separate run of \OPENAIMODEL{}, which had not seen the work, reviewed the whole manuscript refutation-first, including the memory theorem, compared it with the source derivations, and re-ran the corner-lemma programs. It found no false step in the main results, one omitted case in item~(c) of \cref{pf:robust-scope} and one credit sentence, both corrected in this version. It also tested the forced-word lemma (\cref{thr:word}) for $n=1,2$ with two independently written programs.

Version~1.4.0 adds the following. The observation that the forgetting bound holds for all starts at once (\cref{fgt:coalesce}) and its application to all Gibbs measures of the space--time interaction (\cref{cor:gibbs}), through the prefix-cap theorem of \cite{DM26}, were derived by \OPENAIMODEL{} and checked cross-family, refutation-first, by Claude Opus~5.5; a separate run of the same model attacked the load-bearing steps without seeing that check. Neither found a false step. The extension in \cref{pf:gibbs-other} was derived during that check by Claude Opus~5.5 and derived again, in writing this version, by the same model. The stacked automata of \cref{pf:sec-stack} were derived by \OPENAIMODEL{} and checked refutation-first by Claude Opus~5.5. That check found one clause stated without proof, the bound on the mixing time in \cref{cor:stack}(i), and supplied the proof used here; it also proposed the last claim of \cref{cor:stack}(ii)(c), item (ii)(e) and the last claim of (iii), which were derived again in writing this version. The halting case of \cref{cor:stack}(iii) was derived in writing this version. The Lean proofs of \cref{app:lean} were written by agents running Claude Opus~5.5, directed by a lead agent running the same model, which also read the challenge statement of the comparator check against \cref{sec:setting} and \cref{cor:eight}; the formalization led to \cref{app:machine} and to the remark after \cref{thr:corner-plane}. A run of \OPENAIMODEL{} ported the proofs to the toolchain pinned for release, built the comparator challenge, and ran the comparator, the two kernel checkers and the axiom audit. A separate run of \OPENAIMODEL{}, which had not seen the work, then checked the new material of the written version~1.4.0 refutation-first, including \cref{pf:gibbs-other}, the halting case of \cref{cor:stack}(iii), \cref{app:machine} and the fidelity of the formal statement described in \cref{app:lean}, which it read in the Lean sources without building them. It found no failure in the new arguments. It found two statements to correct: \cref{app:lean} claimed the classification for the family of \cref{def:family} itself, which the development proves for an encoded realization once the convention for encoded machines stated there is added; and one sentence of \cref{app:machine} misstated which moves of $M$ take a single step of $M^\circ$. It also asked for corrections of the wording of the one-symbol case in the proof of \cref{cor:gibbs} and of the last clause of the abstract, for the removal of a sentence about earlier work, and for explicit notes where it could not confirm the evidence. All of this is done in this version. Claude Opus~5.5 wrote this version. Models are not listed as authors; the authors take responsibility for every claim.

% =====================================================================
\section*{Verification status}
\addcontentsline{toc}{section}{Verification status}
% =====================================================================

\begin{itemize}
\item \textbf{Proved in the text.} Theorems A, B and C and every lemma and corollary of \cref{sec:setting,sec:forgetting,sec:memory,sec:scaffold,sec:family,sec:refutation,sec:proofs,sec:three} and of \cref{app:machine}, given the cited inputs below; in \cref{sec:three} the one exception is the corner lemma \cref{thr:corner}, which is checked by machine.
\item \textbf{Imported with stated scope.} From \cite{DM26}, where they are proved: Lemma~2.4, Corollary~10.4 and Theorem~10.10(i), that is, the correspondence of \cite{GKLM89} between translation-invariant Gibbs measures of the space--time interaction and stationary trajectories, and the uniqueness of the periodic Gibbs measure of an ergodic positive-rate automaton, used in \cref{cor:gibbs,cor:stack}; the prefix-cap theorem and its corollary for basic flows, Theorem~13.3 and Corollary~13.7, used in \cref{cor:gibbs}(i), \cref{pf:gibbs-other} and \cref{cor:stack}(iii); and, for the remark on the arithmetical degree, Remark~10.9 with Halawi's theorem \cite{Halawi22}.
\item \textbf{Checked less thoroughly.} \Cref{pf:gibbs-other} and the halting case of \cref{cor:stack}(iii) were derived by models of one developer and checked by a model of the other developer only in the review of the written version~1.4.0 (see ``How this work was done'').
\item \textbf{Existence only.} The parser constant $B$ of \cref{thm:parser} is shown to exist and was not computed. Theorem~A's explicit constants come from \cref{cor:eight}, which does not use it.
\item \textbf{Checked by machine.} One finite statement, the corner lemma \cref{thr:corner}, is a computation: an arc-consistency deletion on a box of side $33$ for four roots, with four positive controls. Two independently written programs carry it out and agree, and a third, written separately during review, agrees with them. The programs, their outputs and run receipts are in \texttt{verification/corner-lemma/} of the repository. The review of version~1.3.0 re-ran the programs; the review of version~1.4.0 did not. The third program also finds that the list at the centre empties already on the box of side $13$, the statement used in the remark after \cref{thr:corner-plane}. No other proof uses a computation. The forced-word lemma (\cref{thr:word}), which is proved in the text, was also tested for $n=1$ and $n=2$ by two further independently written programs, for all sixteen pairs of corner letters, with legal-pattern and ablation controls (\texttt{verification/forced-word/}). These tests support the proof and are not part of it.
\item \textbf{Formalized.} \Cref{cor:eight}, Theorem~A on $\Z^8$, and the undecidability of ergodicity in the family are proved in Lean~4 with Mathlib, with no unproved step and only the three standard axioms (\cref{app:lean}). The corner lemma enters there in the form of the remark after \cref{thr:corner-plane}, checked by the Lean kernel. \Cref{thm:width}, \cref{thm:main}, the other corollaries of \cref{sec:proofs} and \cref{fgt:coalesce} are not formalized. The Lean sources are in \texttt{lean/} of the repository. A comparator check confirms that the proved statements are those of a challenge file that imports only Mathlib, and two kernel checkers re-check the proofs (\cref{app:lean}).
\item \textbf{Citations.} Of the references added in version~1.4.0, \cite{Sla81} was read only in its abstract, and \cite{Dob68,DS85,LMS90,vdBS99} are cited for statements recorded in \cite{DM26}; we have not checked them against their full texts.
\item \textbf{Internal only.} Every review is internal; none of it is journal peer review.
\end{itemize}
